\section{工程落地路线图：从研究原型到生产系统}

\subsection{阶段一：可复现实验平台}
先构建可重放环境、统一动作协议、标准评测脚本。阶段目标不是追求最高分，而是建立可诊断、可回归、可扩展的实验底座。

\subsection{阶段二：数据飞轮}
以真实失败任务为核心，持续回流到数据构造与训练流程。教程驱动与合成数据并行推进，但始终保持执行验证约束，防止数据分布漂移。

\subsection{阶段三：生产守护}
加入权限控制、操作审计、异常回滚、人工接管机制。只有当 ``可观测性 + 可恢复性'' 具备，Agent 才能从 demo 进入业务关键链路。

\begin{knowledgebox}{组织层面的协同要求}
多模态 Agent 不是单模型团队可以独立完成的项目。需要环境工程、评估工程、数据工程、模型训练和产品安全团队协作，形成统一发布节奏。
\end{knowledgebox}

\begin{importantbox}{一句话落地原则}
\textit{``先把失败过程看清楚，再去追求更高分数。''} 这比盲目堆参数或堆样本更能提升真实生产力。
\end{importantbox}

\subsection{本章小结}
落地路径应遵循 ``平台先行、数据飞轮、生产守护''。课程的核心价值在于给出了一条可执行、可扩展、可审计的 Agent 工程路线。

